%=============================================================================
% WORK HEALTH AND SAFETY REPORT
%-----------------------------------------------------------------------------
% Drop-in for the \appendices block (IEEEtran): place after Case Study 2
% (order: Case Study 1 -> Case Study 2 -> this).
% Self-contained: no placeholder \ref keys and no bibliography entries; the
% legal instruments are named in text. Requires booktabs (already in use if
% the draft's existing tables compile) and the table* float for full width.
% Risk ratings are qualitative bands (Low / Medium / High). If you prefer to
% carry over the numeric scores from the proposal's project risk register,
% append them in parentheses after each band and keep the matrix consistent.
% Column widths in the register are a starting point; adjust to taste.
%=============================================================================

\section{Work Health and Safety Report}
\label{app:whs}

\subsection{Scope and duty holders}
The placement is screen-based software development and research, with no
laboratory, workshop, fieldwork, chemical or manual-handling exposure. Work
was performed across three settings. Approximately 60 percent of hours were
worked at a home workstation equipped with a laptop, external monitor,
ergonomic chair and sit-stand desk. Approximately 20 percent were worked at
the industry partner's office within the Cicada Westmead health-technology
hub, using a laptop and ergonomic chair, with attendance discretionary
under a remote-first arrangement. The remaining 20 percent were spent at
university facilities, using a laptop in library and meeting-room settings.
Under the \textit{Work Health and Safety Act 2011} (NSW), a student gaining
work experience is a worker (s~7), so the industry partner, the university
and the facility operator hold concurrent duties as persons conducting a
business or undertaking (s~19) and are required to consult, cooperate and
coordinate with one another in respect of the same worker (s~46). The
author holds the worker duties of s~28, taking reasonable care for their
own health and safety and complying with reasonable instructions and
policies. The site-specific expression of these duties is the Cicada
Westmead induction package, completed and signed before office attendance
commenced.

\subsection{Method}
Hazards were identified by reviewing the placement's tasks and settings
against the categories of the model \textit{Code of Practice: How to Manage
Work Health and Safety Risks}, classified by hazard type, and rated
qualitatively by likelihood and consequence before and after controls,
mirroring the project risk register maintained since the proposal.
Psychosocial hazards are classified and controlled with reference to the
psychosocial provisions inserted into the \textit{Work Health and Safety
Regulation 2017} (NSW) in 2022 and the NSW \textit{Code of Practice:
Managing Psychosocial Hazards at Work}. Controls are mapped to the
hierarchy of control measures (reg~36). For hazards inherent to sedentary
screen-based work, elimination and substitution are not reasonably
practicable, so controls concentrate at the engineering and administrative
levels, which the register makes explicit. Table~\ref{tab:whs-register}
presents the resulting register.

\begin{table*}[t]
\centering
\footnotesize
\setlength{\tabcolsep}{3pt}
\caption{Work health and safety risk register for the placement. Ratings
are qualitative (likelihood $\times$ consequence). Controls are tagged with
their hierarchy-of-controls level, where E denotes engineering controls and
A denotes administrative controls.}
\label{tab:whs-register}
\begin{tabular}{p{22mm} p{18mm} p{12mm} p{14mm} p{13mm} p{50mm} p{19mm} p{13mm}}
\toprule
Hazard & Classification & Likelihood & Consequence & Inherent risk &
Controls (hierarchy level) & Duty holders & Residual risk \\
\midrule
Sedentary work and static posture over long sessions & Ergonomic
(musculoskeletal) & Possible & Moderate & Medium & Sit-stand desk enabling
posture alternation at the primary workstation and ergonomic task chairs at
home and the partner office (E). Scheduled posture breaks and rotation
between reading, writing and coding tasks (A). & Worker, partner & Low \\
\addlinespace
Laptop-only working at the partner office and university (about 40 percent
of hours) & Ergonomic (musculoskeletal) & Possible & Moderate & Medium &
The majority of hours run at the fully equipped home workstation, and
sessions at the other sites are shorter and broken by meetings and travel,
limiting exposure (A). Residual exposure is addressed by
recommendation~1. & Worker, university, partner & Low--Medium \\
\addlinespace
Prolonged screen exposure and visual fatigue & Ergonomic (visual) & Likely
& Minor & Medium & External monitor positioned at eye height with adequate
ambient lighting at the primary workstation (E). Regular screen breaks and
batching of off-screen tasks such as reading and meetings (A). & Worker &
Low \\
\addlinespace
High job demands from concurrent thesis and production deadlines,
professional isolation under remote-first work, and work--life boundary
erosion from a workstation within a personal living space & Psychosocial &
Possible & Moderate & Medium & Structured weekly stand-ups with the
partner team and regular academic-supervisor meetings (A). Discretionary
office attendance used for in-person contact (A). Defined working hours
separating work from personal time (A). & Partner, university, worker &
Low \\
\addlinespace
Electrical and trip hazards at the home workstation (trailing leads,
power-board loading) & Physical/electrical & Unlikely & Moderate & Low &
Cable management routing leads clear of walkways, no daisy-chained power
boards, and periodic visual inspection of leads and plugs (E/A). & Worker &
Low \\
\addlinespace
Site-specific hazards of a shared facility (emergency procedures, shared
spaces) & Site/facility & Unlikely & Moderate & Low & Facility induction
completed with signed acknowledgement and compliance with the facility
operator's emergency and house procedures (A). Hazards are managed under
the facility operator's WHS system. & Facility operator, worker & Low \\
\bottomrule
\end{tabular}
\end{table*}

\subsection{Review and recommendations}
No notifiable incidents within the meaning of the Act occurred during the
placement, and no injuries or near-misses of any kind were recorded. The
register is treated as a living document, reviewed when tasks or sites
change and standing for final review at the completion of the placement.
Two improvements are recommended. First, the residual ergonomic exposure
of laptop-only work at the partner office and the university should be
closed with portable peripherals such as a laptop stand and external
keyboard and mouse, or with a pre-use workstation checklist for those
sites. Second, the home workstation controls, currently effective but
informal, should be formalised through a working-from-home workstation
self-assessment agreed with the industry partner, in line with SafeWork
NSW guidance on working from home, so that the arrangement under which the
majority of placement hours were worked is documented rather than assumed.